\documentclass[11pt]{article}

\usepackage[a4paper,margin=1in]{geometry}
\usepackage{amsmath,amssymb,bm}
\usepackage{booktabs}
\usepackage{enumitem}
\usepackage{xcolor}
\usepackage{hyperref}
\usepackage{listings}

\hypersetup{
  colorlinks=true,
  linkcolor=blue,
  urlcolor=blue,
  citecolor=blue
}

\lstset{
  basicstyle=\ttfamily\small,
  breaklines=true,
  frame=single,
  columns=fullflexible
}

\newcommand{\E}{\bm{E}}
\newcommand{\F}{\bm{F}}
\newcommand{\C}{\bm{C}}
\newcommand{\Sstress}{\bm{S}}
\newcommand{\I}{\bm{I}}
\newcommand{\X}{\bm{X}}
\newcommand{\x}{\bm{x}}
\newcommand{\uvec}{\bm{u}}
\newcommand{\R}{\mathbb{R}}
\newcommand{\dd}{\,\mathrm{d}}

\title{Beginner Explanation of the ICKAN Test on the Neo-Hookean RVE Database}
\author{Sebastian ICKAN Tests}
\date{\today}

\begin{document}

\maketitle

\tableofcontents

\section{Purpose of this note}

This note explains, from the beginning, what the ICKAN test is trying to do with
the RVE database generated in this repository. It assumes that the reader is new
to:

\begin{itemize}[leftmargin=2em]
  \item finite strain mechanics,
  \item hyperelasticity,
  \item homogenization,
  \item principal stretches,
  \item invariants,
  \item convexity, monotonicity, and polyconvexity,
  \item KAN/ICKAN-type neural networks.
\end{itemize}

The goal is not only to explain the code, but also to explain the mechanics
ideas behind the code. The main message is:

\begin{quote}
The ICKAN scripts try to learn an effective hyperelastic constitutive law of the
perforated RVE. The input is a macroscopic strain state. The output is the
macroscopic homogenized stress. Instead of predicting stress directly, the ICKAN
predicts a scalar energy, and stress is obtained by differentiating that energy.
\end{quote}

\section{Very short summary}

The Stage 1 FOM database contains pairs of the form
\[
  \text{strain} \longmapsto \text{stress}.
\]
More specifically, every sample looks like
\[
  \E =
  \begin{bmatrix}
    E_{xx} \\
    E_{yy} \\
    G_{xy}
  \end{bmatrix}
  \quad \longmapsto \quad
  \Sstress =
  \begin{bmatrix}
    S_{xx} \\
    S_{yy} \\
    S_{xy}
  \end{bmatrix}.
\]
Here:

\begin{itemize}[leftmargin=2em]
  \item $\E$ is the Green--Lagrange strain,
  \item $G_{xy}$ is the engineering shear strain, so $G_{xy}=2E_{xy}$,
  \item $\Sstress$ is the second Piola--Kirchhoff stress,
  \item both are homogenized, meaning they are average macroscopic quantities
        extracted from the detailed RVE finite element solution.
\end{itemize}

The ICKAN code tries to learn a scalar function
\[
  W = W(\E),
\]
where $W$ is an effective strain energy density. Then stress is computed as
\[
  \Sstress = \frac{\partial W}{\partial \E}.
\]
This is the key idea. The neural network does not directly output three stress
components. It outputs one scalar energy. The stress is obtained from the
gradient of this energy with respect to strain.

\section{Files involved}

The relevant files are:

\begin{itemize}[leftmargin=2em]
  \item \texttt{stage\_1\_training\_set\_fom/trajectory\_i/trajectory\_i\_strain.npy}
  \item \texttt{stage\_1\_training\_set\_fom/trajectory\_i/trajectory\_i\_stress.npy}
  \item \texttt{stage\_1\_training\_set\_fom/trajectory\_i/trajectory\_i\_applied\_strain.npy}
  \item \texttt{ICKAN\_surrogate.py}
  \item \texttt{run\_ICKAN\_surrogate.py}
\end{itemize}

The ICKAN scripts currently load \texttt{trajectory\_i\_strain.npy} and
\texttt{trajectory\_i\_stress.npy}. They do not use the displacement field
\texttt{trajectory\_i\_U.npy}.

\section{The physical problem behind the data}

\subsection{What is the RVE?}

RVE means representative volume element. In this repository, the RVE is a small
two-dimensional sample of material:

\begin{itemize}[leftmargin=2em]
  \item the outer domain is rectangular/square,
  \item there is a circular hole in the center,
  \item the mesh is fixed,
  \item the material is hyperelastic and plane strain,
  \item different macroscopic strain states are imposed on the boundary.
\end{itemize}

The geometry is not being varied in the ICKAN test. The hole size, mesh, and
material parameters are fixed. What changes from sample to sample is the imposed
macroscopic strain.

\subsection{Why do we need homogenization?}

The RVE has local stress and strain fields that vary from point to point. Near
the hole, stress concentrations appear. Farther from the hole, the response is
different.

However, for a macroscopic material model, we do not want to store every local
finite element value. We want an average effective law:
\[
  \text{macroscopic strain} \longmapsto \text{macroscopic stress}.
\]

Homogenization turns the detailed finite element result into average quantities.
The resulting effective material is not simply the original Neo-Hookean material
because the hole changes the macroscopic behavior. The ICKAN is trying to learn
this effective behavior.

\section{Basic mechanics vocabulary}

\subsection{Reference and current configuration}

In finite deformation mechanics, we distinguish between:

\begin{itemize}[leftmargin=2em]
  \item the reference configuration: the undeformed geometry,
  \item the current configuration: the deformed geometry.
\end{itemize}

A material point starts at reference position
\[
  \X =
  \begin{bmatrix}
    X \\ Y
  \end{bmatrix}.
\]
After deformation, it moves to
\[
  \x =
  \begin{bmatrix}
    x \\ y
  \end{bmatrix}.
\]
The displacement is
\[
  \uvec = \x - \X.
\]

\subsection{Deformation gradient}

The deformation gradient is
\[
  \F = \frac{\partial \x}{\partial \X}.
\]
It tells us how an infinitesimal material vector changes during deformation.

For small deformation, one often talks directly about displacement gradients.
For large deformation, $\F$ is more fundamental.

\subsection{Right Cauchy--Green tensor}

From $\F$, we define
\[
  \C = \F^T \F.
\]
This is called the right Cauchy--Green tensor. It measures stretching while
removing rigid-body rotation.

Why does $\C$ remove rotations? Suppose the body only rotates and does not
stretch. Then
\[
  \F = \bm{R},
\]
where $\bm{R}$ is a rotation matrix. A rotation satisfies
\[
  \bm{R}^T\bm{R} = \I.
\]
Therefore
\[
  \C = \F^T\F = \bm{R}^T\bm{R} = \I.
\]
So $\C$ says ``no stretch'', which is correct.

\subsection{Green--Lagrange strain}

The Green--Lagrange strain is
\[
  \E = \frac{1}{2}(\C-\I).
\]
Equivalently,
\[
  \C = 2\E + \I.
\]

This is the strain measure used in your data. In two dimensions:
\[
  \E =
  \begin{bmatrix}
    E_{xx} & E_{xy} \\
    E_{xy} & E_{yy}
  \end{bmatrix}.
\]

The code stores strain in Voigt form:
\[
  \E_{\mathrm{voigt}}
  =
  \begin{bmatrix}
    E_{xx} \\
    E_{yy} \\
    G_{xy}
  \end{bmatrix},
  \qquad
  G_{xy}=2E_{xy}.
\]
Therefore, when the code reconstructs the tensor, it does:
\[
  E_{xy} = \frac{1}{2}G_{xy}.
\]

\section{From strain to boundary displacement}

\subsection{Why the solver needs the deformation gradient F}

The finite element solver imposes displacements on the boundary. But the dataset
is defined in terms of strain $\E$. So the solver must convert $\E$ into a
deformation gradient $\F$.

The code uses:
\[
  \C = 2\E + \I,
\]
then
\[
  \F = \C^{1/2}.
\]
Here $\C^{1/2}$ means the symmetric positive definite square root of $\C$.

\subsection{Why choose F equal to the square root of C?}

Many different deformation gradients can have the same $\C$. For example:
\[
  \F = \bm{R}\C^{1/2}
\]
has the same $\C$ for any rotation $\bm{R}$, because
\[
  \F^T\F
  =
  (\bm{R}\C^{1/2})^T(\bm{R}\C^{1/2})
  =
  \C^{1/2}\bm{R}^T\bm{R}\C^{1/2}
  =
  \C.
\]

The code chooses the rotation-free representative:
\[
  \F = \C^{1/2}.
\]
This is a good choice for machine learning because rigid rotations should not
affect the material law. If rotations were included arbitrarily, the same strain
could correspond to many different boundary displacements, making the learning
problem messier.

\subsection{Boundary displacement}

The imposed affine displacement is
\[
  \uvec_{\mathrm{aff}} = (\F-\I)(\X-\X_c),
\]
where $\X_c$ is the center of the RVE domain. This means the boundary is moved in
a way consistent with the desired macroscopic strain.

\section{Stress measure}

\subsection{What stress is stored?}

The stored stress is the second Piola--Kirchhoff stress:
\[
  \Sstress.
\]

This stress is naturally paired with the Green--Lagrange strain in energy
calculations:
\[
  \dd W = \Sstress : \dd \E.
\]

The double dot product means:
\[
  \Sstress : \dd \E
  =
  S_{xx}\dd E_{xx}
  +
  S_{yy}\dd E_{yy}
  +
  2S_{xy}\dd E_{xy}.
\]

Since the code stores $G_{xy}=2E_{xy}$, the work expression in Voigt form becomes
\[
  \dd W
  =
  S_{xx}\dd E_{xx}
  +
  S_{yy}\dd E_{yy}
  +
  S_{xy}\dd G_{xy}.
\]

This explains why the code can use the simple dot product:
\[
  \Delta W \approx \Sstress_{\mathrm{voigt}}\cdot \Delta \E_{\mathrm{voigt}}.
\]

\section{Principal stretches}

\subsection{What is a principal direction?}

Imagine stretching a square rubber sheet. In some directions the material may be
stretched more, and in other directions less. The special directions where the
stretch happens without shear are called principal directions.

The stretch factors in those directions are called principal stretches:
\[
  \lambda_1,\lambda_2,\lambda_3.
\]

If $\lambda_i > 1$, the material is stretched in that direction. If
$\lambda_i < 1$, it is compressed in that direction.

\subsection{How do principal stretches come from C?}

The eigenvalues of $\C$ are the squares of the principal stretches:
\[
  \text{eigenvalues of } \C
  =
  \lambda_i^2.
\]
Therefore:
\[
  \lambda_i = \sqrt{\text{eigenvalue}_i(\C)}.
\]

The ICKAN code does this:

\begin{lstlisting}[language=Python]
square_eigenvalues = torch.linalg.eigvalsh(C)
eigenvalues = torch.sqrt(square_eigenvalues)
\end{lstlisting}

Despite the variable name \texttt{eigenvalues}, this quantity is actually the
principal stretches $\lambda_i$.

\section{Invariants}

\subsection{Beginner definition}

An invariant is a quantity that does not change when we rotate the coordinate
system.

For example, suppose you describe a deformation in the $x$-$y$ axes, then rotate
your axes and describe the same deformation again. The tensor components change,
but the physical deformation is the same. A good material model should not
change its prediction just because we changed the coordinate system.

Quantities such as:
\[
  \lambda_1,\lambda_2,\lambda_3
\]
or combinations such as:
\[
  I_1 = \mathrm{tr}(\C),
  \qquad
  I_2 = \frac{1}{2}\left[(\mathrm{tr}\C)^2-\mathrm{tr}(\C^2)\right],
  \qquad
  I_3 = \det(\C)
\]
are invariants.

\subsection{Why invariants matter for machine learning}

If a neural network learns directly from tensor components, it may learn a
coordinate-dependent representation. That can be inefficient and physically
fragile.

If the network learns from invariants or principal stretches, it is encouraged
to learn a material response that depends on physical deformation, not on the
chosen coordinate axes.

\subsection{What invariants are used in the current ICKAN code?}

The code does not use $I_1,I_2,I_3$ explicitly. Instead, it uses stretch-like
features:
\[
  \bar{\lambda}_1 = \lambda_1 J^{-1/3},
  \qquad
  \bar{\lambda}_2 = \lambda_2 J^{-1/3},
\]
and a volumetric feature:
\[
  (J-1)^2.
\]

Here
\[
  J = \sqrt{\det(\C)}.
\]

In three-dimensional finite strain mechanics, $J=\det(\F)$ is the volume change.
If $J=1$, the deformation preserves volume. If $J>1$, the material locally
expands. If $J<1$, it locally compresses.

\section{Isochoric and volumetric split}

\subsection{What does isochoric mean?}

Isochoric means volume-preserving. It describes the part of the deformation that
changes shape but not volume.

Volumetric means related to volume change.

Many hyperelastic models split the energy into two parts:
\[
  W = W_{\mathrm{iso}} + W_{\mathrm{vol}}.
\]

The isochoric part controls distortional/shear-like deformation. The volumetric
part controls compression or expansion.

\subsection{Why multiply by J to the power minus one third?}

In three dimensions, the product of the three principal stretches is:
\[
  J = \lambda_1\lambda_2\lambda_3.
\]

The modified stretches
\[
  \bar{\lambda}_i = J^{-1/3}\lambda_i
\]
have product:
\[
  \bar{\lambda}_1\bar{\lambda}_2\bar{\lambda}_3
  =
  J^{-1}\lambda_1\lambda_2\lambda_3
  =
  1.
\]

So the modified stretches remove the volume change and keep only the shape
change.

\subsection{Important plane-strain subtlety}

Your problem is two-dimensional plane strain, but the constitutive law is still
a finite-strain solid law. In plane strain, one can think of an out-of-plane
stretch $\lambda_3=1$. Then:
\[
  J = \lambda_1\lambda_2\lambda_3
  =
  \lambda_1\lambda_2.
\]

The code uses a two-dimensional $\C$ and computes:
\[
  J = \sqrt{\det(\C_{2D})}.
\]
This is compatible with the idea $\lambda_3=1$, but it is worth remembering that
the ICKAN input only includes two in-plane modified stretches plus one
volumetric term.

\section{Hyperelasticity}

\subsection{What is a hyperelastic material?}

A hyperelastic material is a material where stress comes from an energy
function. For Green--Lagrange strain and second Piola--Kirchhoff stress:
\[
  \Sstress = \frac{\partial W}{\partial \E}.
\]

This is similar to a spring. For a linear spring:
\[
  W(u)=\frac{1}{2}ku^2,
  \qquad
  f=\frac{\dd W}{\dd u}=ku.
\]
Energy gives force by differentiation.

For a hyperelastic solid:
\[
  W(\E) \quad \Rightarrow \quad \Sstress(\E).
\]

\subsection{Why learn energy instead of stress directly?}

There are three main reasons.

\paragraph{Reason 1: mechanical consistency.}
If stress is the gradient of an energy, the model is automatically conservative.
That means the stress comes from a potential.

\paragraph{Reason 2: symmetric tangent stiffness.}
The material tangent is the derivative of stress with respect to strain:
\[
  \mathbb{C}_{\mathrm{tan}}
  =
  \frac{\partial \Sstress}{\partial \E}
  =
  \frac{\partial^2 W}{\partial \E^2}.
\]
Second derivatives of a smooth scalar potential are symmetric. This is a useful
physical and numerical property.

\paragraph{Reason 3: better extrapolation structure.}
A scalar energy model can encode physical structure more naturally than three
independent stress outputs.

\subsection{What does the current ICKAN model learn?}

The model learns:
\[
  W_{\theta}(\E),
\]
where $\theta$ denotes the trainable KAN parameters. Then it computes:
\[
  \Sstress_{\theta}(\E)
  =
  \frac{\partial W_{\theta}}{\partial \E}
\]
using PyTorch automatic differentiation.

\section{What is a KAN?}

\subsection{Standard neural network intuition}

A standard multilayer perceptron uses linear combinations of inputs followed by
activation functions. Schematically:
\[
  \bm{z}_{\ell+1}
  =
  \sigma(\bm{W}_{\ell}\bm{z}_{\ell}+\bm{b}_{\ell}).
\]

The trainable objects are mostly the matrix weights $\bm{W}_{\ell}$ and biases
$\bm{b}_{\ell}$.

\subsection{KAN intuition}

KAN stands for Kolmogorov--Arnold Network. In a KAN, the trainable functions are
often placed on edges of the network. Instead of just multiplying by a scalar
weight, an edge can apply a learnable one-dimensional function.

Very roughly, a KAN layer can be imagined as:
\[
  z_j^{\ell+1}
  =
  \sum_i \phi_{ij}^{\ell}(z_i^{\ell}),
\]
where each $\phi_{ij}^{\ell}$ is a learnable scalar function.

In many KAN implementations, these scalar functions are represented by splines.
That is why the code talks about:

\begin{itemize}[leftmargin=2em]
  \item grid size,
  \item spline degree $k$,
  \item grid update from samples.
\end{itemize}

\subsection{What does \texttt{grid\_size} mean?}

The spline functions live on a grid. A larger grid gives more local flexibility,
but can overfit and increases cost.

In the current code:
\[
  \texttt{grid\_size}=4.
\]
This is quite small, meaning the network is intentionally simple.

\subsection{What does spline degree k mean?}

The code uses:
\[
  k=3.
\]
This means cubic splines. Cubic splines are smooth piecewise-polynomial
functions. Smoothness is important because the model differentiates energy to
obtain stress.

\section{What is ICKAN in this repository?}

The file \texttt{ICKAN\_surrogate.py} defines a class:
\[
  \texttt{ICKAN\_W\_Surrogate}.
\]

The ``W'' in the name is important. This is an energy surrogate.

\subsection{Input size}

The code defines:
\[
  \texttt{input\_size}=2\cdot\texttt{order\_stretches}+1.
\]

With the current setting:
\[
  \texttt{order\_stretches}=1,
\]
so:
\[
  \texttt{input\_size}=3.
\]

The three inputs are:
\[
  \bar{\lambda}_1,\quad
  \bar{\lambda}_2,\quad
  (J-1)^2.
\]

\subsection{Network width}

The current code uses:
\[
  \texttt{W\_width}=[3,2,1,1].
\]
This means the KAN maps:
\[
  3 \longrightarrow 2 \longrightarrow 1 \longrightarrow 1.
\]

The final output is a scalar:
\[
  W_{\theta}.
\]

\subsection{Zero-energy normalization}

The code subtracts the network output at zero strain:
\[
  W_{\theta}(\E) - W_{\theta}(\bm{0}).
\]
This enforces:
\[
  W_{\theta}(\bm{0})=0.
\]

That is physically reasonable because the undeformed state should have zero
stored energy, up to an arbitrary additive constant.

\subsection{Zero-stress normalization}

The code also subtracts the stress at zero strain:
\[
  \Sstress_{\theta}(\E) - \Sstress_{\theta}(\bm{0}).
\]
This enforces:
\[
  \Sstress_{\theta}(\bm{0})=\bm{0}.
\]

That is also physically reasonable for an initially stress-free RVE.

\section{How the training data is built}

\subsection{Loading trajectories}

The script \texttt{run\_ICKAN\_surrogate.py} loads ten trajectories:
\[
  i=1,\dots,10.
\]

For each trajectory it loads:
\[
  \texttt{trajectory\_i\_strain.npy},
  \qquad
  \texttt{trajectory\_i\_stress.npy}.
\]

These files contain arrays of shape:
\[
  N_i \times 3.
\]
The number of steps $N_i$ differs by trajectory.

\subsection{Downsampling}

The script sets:
\[
  \texttt{min\_steps}=500.
\]
Then each trajectory is sampled at 500 equally spaced indices.

So the training set becomes:
\[
  10 \times 500 = 5000
\]
samples.

\subsection{Normalization}

The script normalizes stress:
\[
  \Sstress_{\mathrm{norm}}
  =
  \frac{\Sstress}{\max |\Sstress|}.
\]
It also normalizes strain:
\[
  \E_{\mathrm{norm}}
  =
  \frac{\E}{\max |\E|}.
\]

This is common in neural network training because it keeps values near order
one. However, there is an important mechanical caveat:

\begin{quote}
After normalization, the code computes $\C=2\E+\I$ using the normalized strain.
That means the principal stretches and $J$ are computed from normalized strain,
not from physical strain.
\end{quote}

For pure regression, this can still work. But if we want the features
$\lambda_i$ and $J$ to have their physical meaning, then the feature calculation
should be done before normalization, or the normalization should be handled very
carefully.

\subsection{Energy reconstruction}

The FOM database stores strain and stress, but not energy. The script therefore
constructs an approximate energy by numerical integration:
\[
  \Delta W_i
  \approx
  \frac{1}{2}
  \left(
    \Sstress_i+\Sstress_{i-1}
  \right)
  :
  \left(
    \E_i-\E_{i-1}
  \right).
\]

Then:
\[
  W_i = \sum_{j=1}^i \Delta W_j.
\]

This is the trapezoidal rule. It is more accurate than using only the stress at
the beginning or end of each increment.

\subsection{Important path-dependence warning}

For a perfectly hyperelastic material, the energy should be path-independent:
the energy at a strain state should depend only on the final strain, not on the
path used to reach it.

However, the script constructs $W$ by integrating along each sampled trajectory.
If two trajectories reach the same strain through different paths, numerical
integration error, homogenization drift, or data inconsistency could produce
slightly different energies.

The current training mainly uses stress loss, so this is less severe than if it
trained only on $W$. Still, it is a useful thing to remember.

\section{Training objective}

\subsection{The available training modes}

The function \texttt{TRAIN\_KAN} supports three ideas:

\begin{enumerate}[leftmargin=2em]
  \item train only on stress,
  \item train only on energy,
  \item train on a mixed energy + stress loss.
\end{enumerate}

The current settings are:

\begin{lstlisting}[language=Python]
train_W = False
mixed_sovolev_training = False
\end{lstlisting}

Therefore, the current default is stress-based training:
\[
  \mathcal{L}
  =
  \frac{
    \mathrm{mean}\left[
      \left(
        \Sstress_{\theta}
        -
        \Sstress_{\mathrm{FOM}}
      \right)^2
    \right]
  }{
    \mathrm{mean}\left[
      \Sstress_{\mathrm{FOM}}^2
    \right]
    +
    \varepsilon
  }.
\]

\subsection{Why is this still energy-based?}

Even though the loss compares stresses, the model output is still energy. The
stress used in the loss is:
\[
  \Sstress_{\theta}
  =
  \frac{\partial W_{\theta}}{\partial \E}.
\]

So the model is trained so that the derivative of its energy matches the FOM
stress.

\section{Convexity}

\subsection{Convexity in one dimension}

A scalar function $f(x)$ is convex if the line segment between any two points on
the graph lies above the graph.

An intuitive example:
\[
  f(x)=x^2
\]
is convex.

A non-convex example:
\[
  f(x)=\sin(x)
\]
is not convex over a large interval because it curves up and down.

For a smooth one-dimensional function, convexity is equivalent to:
\[
  f''(x) \ge 0.
\]
That means the slope $f'(x)$ is nondecreasing.

\subsection{Convexity in several dimensions}

For a smooth function $f(\bm{x})$, convexity means the Hessian matrix is positive
semidefinite:
\[
  \nabla^2 f(\bm{x}) \succeq 0.
\]

This means:
\[
  \bm{v}^T \nabla^2 f(\bm{x}) \bm{v} \ge 0
  \quad
  \text{for every direction } \bm{v}.
\]

\subsection{What would convex energy imply?}

If
\[
  \Sstress = \nabla_{\E}W,
\]
and $W$ is convex in $\E$, then the stress-strain map is monotone:
\[
  \left(
    \Sstress(\E_1)-\Sstress(\E_2)
  \right)
  :
  \left(
    \E_1-\E_2
  \right)
  \ge 0.
\]

In simple terms: increasing strain should not produce a negative tangent energy
response.

\subsection{Is convexity always required?}

This is subtle. In small-strain elasticity, convexity of energy in strain is a
natural stability condition.

In finite-strain elasticity, requiring the energy to be convex directly in $\F$
is usually too strong and can be physically unrealistic. This is why
polyconvexity is often discussed instead.

\section{Monotonicity}

\subsection{One-dimensional monotonicity}

A function $g(x)$ is monotone increasing if:
\[
  x_1 > x_2
  \quad \Rightarrow \quad
  g(x_1) \ge g(x_2).
\]

For stress-strain behavior in one dimension:
\[
  S = S(E),
\]
monotonicity means that increasing strain does not decrease stress.

\subsection{Tangent interpretation}

If $S(E)$ is differentiable, monotonicity means:
\[
  \frac{\dd S}{\dd E} \ge 0.
\]

Since $S=\dd W/\dd E$, this is:
\[
  \frac{\dd^2 W}{\dd E^2} \ge 0.
\]
So monotonicity of stress is connected to convexity of energy.

\subsection{Multidimensional monotonicity}

For tensor-valued stress-strain maps, monotonicity is written as:
\[
  \left(
    \Sstress(\E_1)-\Sstress(\E_2)
  \right)
  :
  \left(
    \E_1-\E_2
  \right)
  \ge 0.
\]

This says: when moving from strain state $\E_2$ to strain state $\E_1$, the
stress change should do nonnegative work along that strain change.

\section{Polyconvexity}

\subsection{Why polyconvexity exists}

Finite-strain elasticity is usually written in terms of the deformation gradient
$\F$. A hyperelastic energy can be written as:
\[
  W = W(\F).
\]

Direct convexity in $\F$ is often too restrictive. It can conflict with basic
physical requirements such as frame invariance.

Polyconvexity is a weaker condition that is still mathematically useful.

\subsection{Beginner definition}

In three dimensions, an energy $W(\F)$ is polyconvex if it can be written as:
\[
  W(\F) = g(\F,\operatorname{cof}\F,\det\F),
\]
where $g$ is convex in all of its arguments.

Here:

\begin{itemize}[leftmargin=2em]
  \item $\F$ describes line deformation,
  \item $\operatorname{cof}\F$ describes area deformation,
  \item $\det\F$ describes volume deformation.
\end{itemize}

The word ``poly'' refers to depending convexly on several deformation measures,
not only on $\F$ itself.

\subsection{Why people care about polyconvexity}

Polyconvexity is important because it helps guarantee:

\begin{itemize}[leftmargin=2em]
  \item existence of minimizers in nonlinear elasticity,
  \item better mathematical well-posedness,
  \item prevention of some nonphysical material instabilities,
  \item more robust extrapolation under large deformation.
\end{itemize}

\subsection{Does the current ICKAN guarantee polyconvexity?}

No clear guarantee is visible from the current script.

The code uses physically motivated inputs:
\[
  \bar{\lambda}_1,\quad
  \bar{\lambda}_2,\quad
  (J-1)^2.
\]
This is good, but it does not automatically prove convexity or polyconvexity.

To claim polyconvexity, the architecture or training constraints would need to
be designed so that the learned energy satisfies a precise mathematical
condition. The current script should be interpreted as an energy-based surrogate,
not as a certified polyconvex surrogate.

\section{How to read the current ICKAN code}

\subsection{\texttt{ICKAN\_surrogate.py}}

This file defines the model.

\paragraph{Step 1: import external ICKAN package.}

\begin{lstlisting}[language=Python]
sys.path.insert(0, r'C:\ICKANs')
import ickan as KAN
\end{lstlisting}

This path is Windows-specific. On Linux, this import fails unless the package
\texttt{ickan} is installed or the path is changed.

\paragraph{Step 2: build KAN for energy.}

\begin{lstlisting}[language=Python]
self.KAN_W = KAN.MultKAN(...)
\end{lstlisting}

This is the neural network that outputs $W$.

\paragraph{Step 3: convert strain to KAN inputs.}

\begin{lstlisting}[language=Python]
E[:, 0, 0] = strain[:, 0]
E[:, 1, 1] = strain[:, 1]
E[:, 0, 1] = 0.5 * strain[:, 2]
E[:, 1, 0] = 0.5 * strain[:, 2]

C = 2.0 * E + torch.eye(2)
J = torch.linalg.det(C)**0.5
square_eigenvalues = torch.linalg.eigvalsh(C)
eigenvalues = torch.sqrt(square_eigenvalues)
\end{lstlisting}

This reconstructs the mechanics features.

\paragraph{Step 4: evaluate energy.}

\begin{lstlisting}[language=Python]
W_raw = self.KAN_W.forward(kan_input)
return W_raw - W0
\end{lstlisting}

The subtraction makes the energy zero at the reference state.

\paragraph{Step 5: get stress by autograd.}

\begin{lstlisting}[language=Python]
predicted_stress = torch.autograd.grad(
    outputs=W,
    inputs=strain_database,
    grad_outputs=torch.ones_like(W),
    create_graph=True
)[0]
\end{lstlisting}

This is the central hyperelastic idea.

\subsection{\texttt{run\_ICKAN\_surrogate.py}}

This file does the experiment.

\paragraph{Step 1: load data.}

It loops over trajectories 1 to 10 and loads:
\[
  \E_i(t),\quad \Sstress_i(t).
\]

\paragraph{Step 2: downsample.}

Every trajectory is reduced to 500 points.

\paragraph{Step 3: normalize.}

It scales strain and stress to order-one values.

\paragraph{Step 4: integrate energy.}

It computes approximate $W$ from stress-strain work.

\paragraph{Step 5: train ICKAN.}

It trains the energy model, using stress error as the default loss.

\paragraph{Step 6: save plots.}

It saves:

\begin{itemize}[leftmargin=2em]
  \item predicted energy versus reference integrated energy,
  \item predicted stress versus FOM stress,
  \item spline/KAN plots,
  \item model weights.
\end{itemize}

\section{Important caveats in the current scripts}

\subsection{Caveat 1: external dependency path}

The file \texttt{ICKAN\_surrogate.py} contains:
\[
  \texttt{C:\textbackslash ICKANs}.
\]
This is not portable. The code will only run where that path exists or where
\texttt{ickan} is installed.

\subsection{Caveat 2: output directory}

The script saves files into:
\[
  \texttt{ICKAN\_predictions/}.
\]
But the script does not explicitly create that folder. It should call something
like:

\begin{lstlisting}[language=Python]
os.makedirs("ICKAN_predictions/splines", exist_ok=True)
\end{lstlisting}

\subsection{Caveat 3: computed strain versus applied strain}

The script uses:
\[
  \texttt{trajectory\_i\_strain.npy}.
\]
This is the homogenized computed strain. The imposed boundary strain is:
\[
  \texttt{trajectory\_i\_applied\_strain.npy}.
\]

Which one should be used depends on the intended final use.

\begin{itemize}[leftmargin=2em]
  \item If the surrogate will be queried by prescribed macro strain, use applied
        strain.
  \item If the surrogate is intended to map homogenized measured strain to
        homogenized stress, use computed strain.
\end{itemize}

For most constitutive surrogate use cases, the cleaner input is probably the
applied macroscopic strain.

\subsection{Caveat 4: normalization and physical features}

The model computes $\C$, $J$, and stretches from normalized strain. This makes
the inputs numerically convenient, but the quantities are no longer the physical
$\C$, $J$, and stretches.

A more physically transparent workflow would be:

\begin{enumerate}[leftmargin=2em]
  \item start from physical strain,
  \item compute physical invariants/stretches,
  \item normalize those features for neural network training,
  \item save the normalization constants.
\end{enumerate}

\subsection{Caveat 5: no train/test split}

The current script trains and plots using the same data. This checks whether the
network can fit the database, but it does not test generalization.

A better test would be:

\begin{itemize}[leftmargin=2em]
  \item train on 8 trajectories,
  \item validate on 1 trajectory,
  \item test on 1 trajectory,
  \item or leave out an entire $G_{xy}$ level.
\end{itemize}

\subsection{Caveat 6: not certified convex or polyconvex}

The model is energy-based, but that alone does not guarantee:

\begin{itemize}[leftmargin=2em]
  \item convexity,
  \item monotonicity,
  \item polyconvexity,
  \item positive tangent stiffness everywhere.
\end{itemize}

Those properties require either special architecture constraints, special loss
terms, or post-training verification.

\section{What would a good ICKAN test report include?}

A useful report should include:

\begin{enumerate}[leftmargin=2em]
  \item which trajectories were used for training,
  \item which trajectories were held out for validation/testing,
  \item whether inputs are applied strain or homogenized computed strain,
  \item strain and stress normalization constants,
  \item whether invariants are computed before or after normalization,
  \item stress relative error per component,
  \item energy error if energy is used,
  \item monotonicity or tangent checks along selected paths,
  \item plots for unseen trajectories,
  \item saved model weights and reproducibility seed.
\end{enumerate}

\section{Recommended next steps}

\subsection{Step 1: make the code portable}

Replace the hard-coded Windows path with a documented installation or a local
relative path. For example:

\begin{lstlisting}[language=Python]
try:
    import ickan as KAN
except ImportError as exc:
    raise ImportError(
        "Could not import 'ickan'. Install it or add its path to PYTHONPATH."
    ) from exc
\end{lstlisting}

\subsection{Step 2: create output directories}

Before saving:

\begin{lstlisting}[language=Python]
os.makedirs("ICKAN_predictions/splines", exist_ok=True)
\end{lstlisting}

\subsection{Step 3: decide the input strain}

For a constitutive law, prefer:
\[
  \texttt{trajectory\_i\_applied\_strain.npy}
\]
unless there is a specific reason to use homogenized computed strain.

\subsection{Step 4: save scalers}

If strain and stress are normalized, save:

\begin{itemize}[leftmargin=2em]
  \item strain scale,
  \item stress scale,
  \item energy scale,
  \item feature scales.
\end{itemize}

Without these, it is hard to use the model later in a physical simulation.

\subsection{Step 5: add a real test}

For example:

\begin{itemize}[leftmargin=2em]
  \item train on trajectories 1--8,
  \item validate on trajectory 9,
  \item test on trajectory 10.
\end{itemize}

Or more physically:

\begin{itemize}[leftmargin=2em]
  \item train on $G_{xy}=0,0.05,-0.05$,
  \item test on $G_{xy}=0.1,-0.1$.
\end{itemize}

\section{Mental model for Sebastian}

The easiest way to think about the entire pipeline is:

\begin{enumerate}[leftmargin=2em]
  \item The FOM solver is an expensive laboratory.
  \item The laboratory deforms the perforated RVE in many ways.
  \item For each deformation, it reports average strain and average stress.
  \item The ICKAN tries to learn the hidden effective material energy of the
        perforated RVE.
  \item Once the energy is learned, stress is obtained by differentiation.
\end{enumerate}

So the ICKAN is not replacing the mesh directly. It is replacing the final
homogenized constitutive relationship:
\[
  \E \mapsto \Sstress.
\]

\section{Glossary}

\begin{description}[leftmargin=2em]
  \item[RVE] Representative volume element. A small sample used to infer
  macroscopic material behavior.

  \item[Homogenization] Averaging a detailed local field to obtain effective
  macroscopic quantities.

  \item[Green--Lagrange strain] A finite-strain measure
  $\E=\frac{1}{2}(\C-\I)$.

  \item[Second Piola--Kirchhoff stress] A stress measure work-conjugate to
  Green--Lagrange strain.

  \item[Work-conjugate] Two quantities are work-conjugate if their product gives
  mechanical work or energy increment. Here $\Sstress:\dd\E$.

  \item[Deformation gradient] Tensor $\F$ that maps reference differential
  vectors to current differential vectors.

  \item[Right Cauchy--Green tensor] Tensor $\C=\F^T\F$ measuring stretch without
  rigid-body rotation.

  \item[Principal stretches] Stretch factors $\lambda_i$ in the principal
  directions of deformation.

  \item[Invariant] A scalar quantity that does not change under coordinate
  rotation.

  \item[Hyperelasticity] A material theory where stress is obtained as the
  derivative of an energy.

  \item[Convexity] A mathematical property meaning the function has no downward
  curvature. For smooth scalar functions, this is related to nonnegative second
  derivative or positive semidefinite Hessian.

  \item[Monotonicity] A property meaning the output does not decrease when the
  input increases, generalized to tensors through nonnegative work of
  increments.

  \item[Polyconvexity] A finite-strain stability condition where the energy is
  convex in $\F$, $\operatorname{cof}\F$, and $\det\F$ through another function.

  \item[KAN] Kolmogorov--Arnold Network. A network using learnable scalar
  functions, often splines, on edges.

  \item[ICKAN] The specific ICKAN library/model used by Alejandro's script. In
  this repository, it is used as an energy-based constitutive surrogate.
\end{description}

\section{Final interpretation}

Alejandro's current ICKAN test is best interpreted as a first prototype for:
\[
  \boxed{
  \text{learning the effective hyperelastic energy of the perforated RVE}
  }
\]
from your FOM-generated homogenized strain-stress data.

The idea is scientifically reasonable because the FOM material is hyperelastic
and the RVE response should have an effective energy-like behavior under the
controlled loading paths. However, the current implementation should not yet be
interpreted as a validated, portable, polyconvex constitutive model. It is an
energy-based regression experiment that needs careful validation.

\end{document}
